% PDF p. 373
\BackSection{习题}

\begin{enumerate}
  \item \textbf{习题 6.1}\par
  证明 $\langle R(v,e_i)v,e_j\rangle=R_{ikjl}v^kv^l$。

  \textbf{提示}\quad
  $\displaystyle\langle R(v,e_i)v,e_j\rangle
    =\sum_{k,l}R_{ikjl}v^kv^l=R_{ikjl}v^kv^l$。
  \item \textbf{习题 6.2}\par
  计算
  \[
    \int_{S^{n-1}}\omega^k\omega^l\,dS
     =\frac{\operatorname{Vol}(S^{n-1})}{n}\delta^{kl},
  \]
  其中 $\operatorname{Vol}(S^{n-1})=\dfrac{\pi^{n/2}}{\Gamma(n/2+1)}r^n$。
  \translatornote{原书沿用了附录 6F 的球体积写法。单位球面面积见式 \eqref{eq:6F5} 的译注。}
  \item \textbf{习题 6.3}\par
  计算 $\displaystyle I_T=\int_0^{2\pi}\sin(\omega t)\,dW_t$。
  \item \textbf{习题 6.4}\par
  求均值回复的奥恩斯坦–乌伦贝克过程。设 $dX_t=-\theta X_t\,dt+\sigma\,dW_t$，求解 $X(t)$。
  \item \textbf{习题 6.5}\par
  取 $f(x)=e^{\lambda x}$，计算奥恩斯坦–乌伦贝克过程的 $\mathcal Lf(x)$。
\end{enumerate}

% PDF p. 374
\BackSection{参考文献}
\begin{bookreferences}
  \item Arvanitidis, G., Hansen, L. K. and Hauberg, S. (2017). Latent space oddity: On the curvature of deep generative models. arXiv:1710.11379.
  \item Arvanitidis, G., Hansen, L. K. and Hauberg, S. (2018). Latent space oddity: On the curvature of deep generative models. \emph{Proceedings of International Conference on Learning Representations (ICLR)}, Vancouver.
  \item Bhattacharya, R. and Bhattacharya, R. (2012). \emph{Nonparametric Inference on Manifolds: With Applications to Shape Spaces}. Cambridge University Press, Cambridge.
  \item Chadebec, C. and Allassonnière, S. (2022). A geometric perspective on variational autoencoders. arXiv:2209.07370.
  \item Chen, N., Ferroni, F., Klushyn, A., Paraschos, A., Bayer, J. and van der Smagt, P. (2018). Fast approximate geodesics for deep generative models. arXiv:1812.08284.
  \item Chen, N., Klushyn, A., Ferroni, F., Bayer, J. and van der Smagt, F. (2020). Learning flat latent manifolds with VAEs. arXiv:2002.04881.
  \item Cheng, X. and Xie, Y. (2024). Kernel two-sample tests for manifold data. \href{https://arxiv.org/abs/2105.03425v4}{\nolinkurl{arXiv:2105.03425v4}}.
  \item Chu, L. and Da, X. (2024). Manifold energy two-sample test. \emph{Electronic Journal of Statistics}. 18: 145–166.
  \item Daniel, O., Vic, P., Leif, E., David, G. and Armin, S. (2013). Nonparametric two-sample tests on homogeneous Riemannian manifolds, Cholesky decompositions and diffusion tensor image analysis. \emph{Journal of Multivariate Analysis}. 119: 163–175.
  \item de Surrel, T., Lotte, F., Chevallier, S. and Yger, F. (2025). Wrapped Gaussian on the manifold of symmetric positive definite matrices. arXiv:2502.01512.
  \item Falorsi, L., de Kano, M., Davidson, T. and Tomczak, J. M. (2018). Explorations on homeomorphic variational auto-encoding. arXiv:1807.04689.
  \item Fletcher, T. and Joshi, S. (2007). Riemannian geometry for the statistical analysis of diffusion tensor data. \emph{Signal Processing}. 87(2): 250–262.
  \item Giovanis, D. G., Crabtree, E., Ghanem, R. G. and Kevrekidis, I. G. (2025). Generative learning of densities on manifolds. arXiv:2503.03963.
  \item Hallou, A., He, R., Simons, R. D. and Dumitrascu, B. (2025). A computational pipeline for spatial mechano-transcriptomics. \emph{Nature Methods}. 22: 737–750. \url{https://doi.org/10.1038/s41592-025-02618-1}
  \item Ho, J., Jain, A. and Abbeel, P. (2020). Denoising diffusion probabilistic models. In Larochelle, H., Ranzato, M., Hadsell, R., Balcan, M. and Lin, H., editors, \emph{Advances in Neural Information Processing Systems} 33: 6840–6851. Annual Conference on Neural Information Processing Systems 2020, NeurIPS 2020, December 6–12, 2020, virtual. 2020.
  \item Kalatzis, D., Eklund, D., Arvanitidis, G. and Hauberg, S. (2020). Variational autoencoders with Riemannian Brownian Motion priors. arXiv:2002.05227.
  \item Kingma, D. P. and Welling, M. (2014). Auto-encoding variational bayes. In \emph{Proceedings of the 2nd International Conference on Learning Representations (ICLR)}, Banff, AB.
  \item Liu, S., Wang, L. and Wang, Y. (2025). Guiding time-varying generative models with natural gradients on exponential family manifold. arXiv:2502.07650.
  \item Miolane, N. and Holmes, S. (2019). Learning weighted submanifolds with variational autoencoders and Riemannian variational autoencoders. arXiv:1911.08147.
  \item Øksendal, B. (2023). \emph{Stochastic Differential Equations: An Introduction with Applications}. Springer Berlin, Heidelberg.
  \item Park, J. Y., Lee, Y., Kim, T.-J. and Choi, J.-H. (2025). Riemannian geometric-based meta learning. arXiv:2503.10993.
  \item Paul's Online Notes. (2025). Section 13.7: Directional Derivatives. \url{https://tutorial.math.lamar.edu/classes/calciii/directionalderiv.aspx}
  \item Randy, S. (2023). Covariant derivatives and curvature. \url{https://cphysics.org/article/03519}
  \item Rozo, L., González-Duque, M., Jaquier, N. and Hauberg, S. (2025). Riemann\textsuperscript{2}: Learning Riemannian submanifolds from Riemannian data. arXiv:2503.05540.
% PDF p. 375
  \item Shamsolmoali, P., Zareapoor, M., Zhou, H., Tao, D. and Li, X. (2023). VTAE: Variational transformer autoencoder with manifolds learning. arXiv:2304.00948.
  \item Wang, J., Chen, M., Zhao, T, Liao, W. and Xie Y. (2022). A manifold two-sample test study: Integral probability metric with neural networks. arXiv:2205.02043.
  \item Whiteley, N., Gray, A. and Rubin-Delanchy, P. (2025). \emph{Statistical exploration of the Manifold Hypothesis}. arXiv:2208.11665.
  \item Williams, D. J., Wang, L., Ying, Q., Liu, S. and Kolar, M. (2024). High-dimensional differential parameter inference in exponential family using time score matching. arXiv:2410.10637.
  \item Wynne, G. and Duncan, A. B. (2022). A kernel two-sample test for functional data. \emph{Journal of Machine Learning Research}. 23: 1–51.
  \item Xavier Pennec. (2004). Probabilities and Statistics on Riemannian Manifolds: A Geometric approach. RR5093, INRIA. 2004. inria–00071490.
\end{bookreferences}
